% !PRINTED 444
\weberpage{444}
%<<CONT
man so verfährt, als ob es sich um die Bestimmung des grössten
gemeinschaftlichen Theilers von $m$ und $n$ handelte.

Aus 8. und 9. ergiebt sich noch ein letzter Satz, der gilt, wenn
$n$ und $n'$ zwei positive ungerade zu $m$ theilerfremde Zahlen sind:

\[
10.\quad \left(\frac{m}{n}\right)\left(\frac{m}{n'}\right) = \left(\frac{m}{nn'}\right).
\]

Seine Richtigkeit ergiebt sich, wenn $m$ ungerade und positiv
ist auf Grund der Congruenz

% [sic] Within §138 the source reuses equation label (7), already printed on p.442 for the alpha formula; preserved diplomatically.
\weq{7}{(n-1)(n'-1) = (nn'-1) - (n-1) - (n'-1) \equiv 0 \ (\mathrm{mod}\ 4),}

wenn man auf beide Seiten von 10. die Formel 9. und dann 8.
anwendet, und wenn $m$ gerade oder negativ ist, aus 4. und 6.

Nach diesen Sätzen kann man $\left(\dfrac{m}{n}\right)$ auf ein Product von
Werthen $\left(\dfrac{q}{p}\right)$ zurückführen, worin $p, q$ Primzahlen sind. Die
Berechnung von $\left(\dfrac{m}{n}\right)$ geschieht aber leichter nach 9. ohne diese
Zerlegung.

Wir wollen nun noch eine andere Bedeutung des Symbols
$\left(\dfrac{m}{p}\right)$ kennen lernen, für den Fall, dass $p$ eine ungerade Primzahl
 ist, durch die ursprünglich das Symbol von \gesperrt{Legendre} eingeführt
 ist.

Wenn wir alle durch $p$ nicht theilbaren ganzen Zahlen $x$
ins Quadrat erheben und die Reste der Division durch $p$ aufsuchen,
so erhalten wir im Ganzen nur $\dfrac{p-1}{2}$ verschiedene
Reste; denn erstens geben die Quadrate congruenter Zahlen dieselben
 Reste, und zweitens geben die Quadrate entgegengesetzter
Zahlen auch dieselben Reste. Wir bekommen also gewiss alle
Reste, wenn wir die $\dfrac{p-1}{2}$ Zahlen $\nu$

\[
(\nu)\qquad 1,\ 2,\ 3\cdots\frac{p-1}{2}
\]

ins Quadrat erheben. Auf der anderen Seite ergeben auch die
Quadrate dieser Zahlen $\nu$ lauter verschiedene Reste; denn es
kann, wenn $\nu$ und $\nu'$ verschiedene von ihnen sind, niemals

\[
\nu^2 - \nu'^2 = (\nu - \nu')(\nu + \nu')
\]
%CONT>>
